\documentclass[11pt]{article}
\usepackage[T1]{fontenc}
\usepackage[utf8]{inputenc}
\usepackage{amsmath,amssymb,mathtools}
\usepackage[a4paper,margin=1in]{geometry}
\usepackage{hyperref}

\title{Guide to the Brenier Formalization}
\author{}
\date{}

\begin{document}
\maketitle

\section*{What this file explains}
This note describes the proof architecture presently formalized in the repository
\texttt{OptimalTransport}. It explains:
\begin{itemize}
\item the overall structure of the argument;
\item the main definitions and theorems;
\item where each step is implemented;
\item the precise form of the current Brenier endpoint.
\end{itemize}

\section*{The target statement}
The classical Brenier theorem for the quadratic cost says that on
\[
  E = \mathbb{R}^n,
\]
if \(\mu,\nu\) are probability measures with finite second moments and
\(\mu \ll \mathrm{vol}\), then for every \(\kappa > 0\) there exists a convex
function \(\phi : E \to \mathbb{R}\) such that
\[
  (\nabla \phi)_{\#}\mu = \nu,
\]
and the graph plan
\[
  (\mathrm{id}, \nabla \phi)_{\#}\mu
\]
is the optimal coupling for the cost
\[
  c_\kappa(x,y) = \kappa \|x-y\|^2.
\]

In the present repository, the final theorem has been packaged in a slightly more explicit form:
the last remaining nontrivial support hypothesis is isolated as
\texttt{Coupling.HasRockafellarBound}. Under that hypothesis, the rest of the Brenier proof is
formalized completely.

\section*{Repository map}
\subsection*{1. Couplings and costs}
\begin{itemize}
\item \texttt{OptimalTransport/Coupling.lean}

Defines:
\[
  \texttt{TransportPlan}\ X\ Y := \mathcal P(X \times Y),
\]
the marginals \(\pi \mapsto \pi_1,\pi_2\), graph plans, and the type
\(\texttt{Coupling}\ \mu\ \nu\).

\item \texttt{OptimalTransport/Cost.lean}

Defines the two transport cost functionals:
\[
  \texttt{transportCostENNReal}(c,\pi)
  = \int c\, d\pi,
\]
for \(c : X \to Y \to \mathbb R_{\ge 0}^{\infty}\), and the real-valued refinement
\[
  \texttt{transportCost}(c,\pi)
  = \int c\, d\pi.
\]
This file also defines:
\[
  \texttt{HasFiniteSecondMoment}(\mu),
  \qquad
  \texttt{quadraticCost}\ \kappa(x,y)=\kappa\|x-y\|^2.
\]
\end{itemize}

\subsection*{2. Existence of optimal plans}
\begin{itemize}
\item \texttt{OptimalTransport/Existence.lean}

This file proves existence of optimal couplings on Polish spaces for lower semicontinuous
\(\mathbb R_{\ge 0}^{\infty}\)-valued costs.

The strategy is the standard Villani strategy:
\begin{enumerate}
\item define the fixed-marginal set \(\Gamma(\mu,\nu)\);
\item prove it is closed;
\item prove it is tight by the estimate
  \[
    \pi((K \times L)^c) \le \mu(K^c) + \nu(L^c);
  \]
\item apply Prokhorov to get compactness;
\item prove lower semicontinuity of the cost by Portmanteau and truncation;
\item minimize on the compact set.
\end{enumerate}

The main theorem is
\[
  \texttt{exists\_optimalCoupling\_of\_isLowerSemicontinuousCost}.
\]
\end{itemize}

\subsection*{3. Cyclical monotonicity and subgradients}
\begin{itemize}
\item \texttt{OptimalTransport/CyclicMonotone.lean}

Defines \(c\)-cyclical monotonicity:
\[
  \texttt{CCyclicallyMonotone}\ c\ \Gamma.
\]
It also proves the algebraic rewriting lemma
\[
  \texttt{cCyclicallyMonotone\_of\_pairing},
\]
which turns costs of the form
\[
  a(x)+b(y)-\kappa\langle x,y\rangle
\]
into the pairing inequality.

\item \texttt{OptimalTransport/Subgradient.lean}

Defines the subgradient relation
\[
  y \in \partial \phi(x)
  \iff
  \forall z,\ \phi(x)+\langle y,z-x\rangle \le \phi(z),
\]
as \texttt{Subgradient \(\phi\) x y}, and its graph
\texttt{SubgradientGraph \(\phi\)}.

This file proves:
\begin{itemize}
\item affine and scaling rules for subgradients;
\item cyclical monotonicity of subgradient graphs;
\item the Brenier-friendly cost specialization
  \[
    \texttt{cCyclicallyMonotone\_subgradientGraph\_of\_pairing}.
  \]
\end{itemize}
\end{itemize}

\subsection*{4. Rockafellar's construction}
\begin{itemize}
\item \texttt{OptimalTransport/Rockafellar.lean}

This file constructs the rooted finite-chain Rockafellar potential. The basic chain functional is
\texttt{rockafellarChainValue}, and the rooted value set is
\[
  \texttt{rockafellarValueSet base }\Gamma\ x.
\]

The potential itself is
\[
  \phi_\Gamma(x)
  :=
  \texttt{rockafellarPotential base }\Gamma\ h_{\mathrm{bdd}}(x)
  =
  \sup \texttt{ rockafellarValueSet base }\Gamma\ x.
\]

The main theorems are:
\begin{itemize}
\item \texttt{subset\_SubgradientGraph\_rockafellarPotential};
\item \texttt{convexOn\_rockafellarPotential}.
\end{itemize}

These show that, once the rooted value sets are known to be bounded above, the support set is
contained in the subgradient graph of a convex function.
\end{itemize}

\subsection*{5. Differentiability of convex potentials}
\begin{itemize}
\item \texttt{OptimalTransport/ConvexAE.lean}

This file proves that convex functions are differentiable almost everywhere in finite-dimensional
real spaces:
\[
  \texttt{ConvexOn.ae\_differentiableAt},
  \qquad
  \texttt{ConvexOn.ae\_hasGradientAt}.
\]

It also proves the uniqueness of the subgradient at differentiability points:
\[
  \texttt{Subgradient.eq\_gradient\_of\_hasGradientAt}.
\]

This is the analytic step that collapses \(\partial \phi(x)\) to \(\{\nabla \phi(x)\}\) almost
everywhere.
\end{itemize}

\subsection*{6. From subgradient supports to maps}
\begin{itemize}
\item \texttt{OptimalTransport/Brenier.lean}

The first part of this file proves the measure-theoretic graph lemmas:
\begin{itemize}
\item \texttt{ae\_mem\_graphOn\_of\_support\_subset};
\item \texttt{TransportPlan.eq\_graph\_of\_ae\_mem\_graphOn};
\item \texttt{Coupling.toTransportPlan\_eq\_graph\_of\_support\_subset}.
\end{itemize}

The Milestone 6 bridge is then:
\begin{itemize}
\item \texttt{TransportPlan.ae\_mem\_graphOn\_gradient\_of\_support\_subset\_subgradientGraph};
\item \texttt{TransportPlan.eq\_graph\_gradient\_of\_support\_subset\_subgradientGraph};
\item \texttt{Coupling.snd\_eq\_map\_gradient\_of\_support\_subset\_subgradientGraph}.
\end{itemize}

These results show that if \(\pi\) is supported in the subgradient graph of a convex function and
its first marginal is absolutely continuous with respect to Haar measure, then
\[
  \pi = (\mathrm{id}, \nabla \phi)_{\#}\mu,
  \qquad
  (\nabla \phi)_{\#}\mu = \nu.
\]
\end{itemize}

\section*{Logical structure of the proof}
The formal proof now has the following shape.

\subsection*{Step A. Start from a coupling}
Let \(\pi \in \Pi(\mu,\nu)\).

\subsection*{Step B. Put the support into a convex subgradient graph}
Assume the support of \(\pi\) satisfies the bounded-root Rockafellar condition. Then
\texttt{Rockafellar.lean} builds a convex potential \(\phi\) such that
\[
  \operatorname{supp}(\pi) \subseteq \operatorname{Graph}(\partial \phi).
\]

\subsection*{Step C. Differentiate almost everywhere}
Since \(\phi\) is convex, \(\phi\) is differentiable \(m\)-almost everywhere. If
\(\mu \ll m\), then \(\phi\) is differentiable \(\mu\)-almost everywhere as well.

\subsection*{Step D. Collapse subgradients to gradients}
At those points of differentiability,
\[
  y \in \partial \phi(x)
  \Longrightarrow
  y = \nabla \phi(x).
\]
Hence \(\pi\) is almost surely supported on the graph of \(\nabla \phi\).

\subsection*{Step E. Upgrade almost-sure graph support to exact graph-plan equality}
Using the graph lemmas from \texttt{Brenier.lean}, one gets
\[
  \pi = (\mathrm{id}, \nabla \phi)_{\#}\mu.
\]
Taking the second marginal yields
\[
  (\nabla \phi)_{\#}\mu = \nu.
\]

\subsection*{Step F. Transfer optimality}
If \(\pi\) was already known to be optimal for the quadratic cost, then the graph plan is the same
transport plan, so it is automatically optimal as well.

\section*{The present final theorem}
The final theorem currently proved in the code is:
\begin{quote}
Let \(E\) be a finite-dimensional real inner product space. Let
\(\pi \in \Pi(\mu,\nu)\) be a coupling. Assume:
\begin{enumerate}
\item \(\mu \ll m\), where \(m\) is an additive Haar measure;
\item \(\pi\) is optimal for the quadratic cost \(\kappa\|x-y\|^2\);
\item \(\pi\) satisfies \texttt{Coupling.HasRockafellarBound}.
\end{enumerate}
Then there exists a convex \(\phi : E \to \mathbb R\) such that
\[
  \pi = (\mathrm{id}, \nabla \phi)_{\#}\mu,
  \qquad
  (\nabla \phi)_{\#}\mu = \nu,
\]
and the graph plan of \(\nabla \phi\) is optimal for the same quadratic cost.
\end{quote}

The precise theorem names are:
\begin{itemize}
\item \texttt{Coupling.exists\_convex\_potential\_eq\_graph\_of\_hasRockafellarBound};
\item \texttt{Coupling.exists\_convex\_potential\_of\_optimal\_quadratic\_of\_hasRockafellarBound};
\item Euclidean specialization:
  \texttt{Coupling.exists\_convex\_potential\_of\_optimal\_quadratic\_of\_hasRockafellarBound\_euclidean}.
\end{itemize}

\section*{What remains to reach classical Brenier}
Only one conceptual gap remains between the present code and the classical textbook theorem.

One still needs a theorem deriving the Rockafellar support hypothesis directly from optimal
transport optimality. Concretely, the missing implication is:
\[
  \text{optimal quadratic coupling}
  \Longrightarrow
  \text{support satisfies the Rockafellar boundedness hypothesis}.
\]

In a more classical presentation, this missing step is usually supplied by:
\begin{enumerate}
\item optimality \(\Rightarrow\) cyclical monotonicity of the support;
\item cyclical monotonicity \(\Rightarrow\) Rockafellar potential with the right boundedness
  properties.
\end{enumerate}

Everything after that implication is now formalized.

\section*{Where to look first}
For a quick code walk-through, the most important theorems are:
\begin{itemize}
\item \texttt{OptimalTransport/Existence.lean}:
  \texttt{exists\_optimalCoupling\_of\_isLowerSemicontinuousCost};
\item \texttt{OptimalTransport/Rockafellar.lean}:
  \texttt{subset\_SubgradientGraph\_rockafellarPotential},
  \texttt{convexOn\_rockafellarPotential};
\item \texttt{OptimalTransport/ConvexAE.lean}:
  \texttt{ConvexOn.ae\_hasGradientAt},
  \texttt{Subgradient.eq\_gradient\_of\_hasGradientAt};
\item \texttt{OptimalTransport/Brenier.lean}:
  \texttt{TransportPlan.eq\_graph\_gradient\_of\_support\_subset\_subgradientGraph},
  \texttt{Coupling.exists\_convex\_potential\_of\_optimal\_quadratic\_of\_hasRockafellarBound}.
\end{itemize}

\end{document}
